\section{Introduction}
\label{sec:intro}

When deploying large language models (LLMs) as assistants and agents, practitioners routinely embed conditional rules in system prompts: ``if the user asks about pricing, respond with the latest rate card''; ``if the input contains profanity, decline to answer.'' These \emph{if-then instructions} govern when and how the model should act, and their reliability directly determines the trustworthiness of the deployed system. Yet when such rules misfire---triggering when they should not or failing to trigger when they should---the consequences range from poor user experience to safety violations.

A growing body of work evaluates LLM instruction following, including IFEval~\citep{zhou2023ifeval}, FollowBench~\citep{jiang2023followbench}, and COLLIE~\citep{yao2023collie}. These benchmarks focus on \emph{imperative} instructions (``write in JSON'', ``keep it under 100 words'') but do not isolate the \emph{conditional} aspect: whether the model correctly determines when a rule applies. The SIFo benchmark~\citep{chen2024sifo} includes a security-rules task that is implicitly conditional, and \citet{heo2024llms} show that internal instruction-following signals do not generalize across instruction types---hinting that conditional capacity may vary systematically. However, no benchmark directly and comprehensively evaluates if-then instruction following.

We fill this gap with \benchname, the first systematic benchmark for conditional instruction following. We design 120 test cases spanning 8 instruction types that cover a range of cognitive demands: simple keyword matching, format switching, content detection, length control, negation, nested branching, conjunction, and disjunction. Each test case pairs a conditional instruction with an input that either triggers or does not trigger the condition, and compliance is verified by deterministic checker functions inspired by IFEval's verifiable-instruction paradigm~\citep{zhou2023ifeval}. We evaluate two models from the GPT family---\gptlarge and \gptsmall---and analyze accuracy, false positive rates, and false negative rates per instruction type.

Our main findings are:
\begin{itemize}[leftmargin=*,itemsep=0pt,topsep=0pt]
    \item We introduce \benchname, a benchmark of 8 conditional instruction types with balanced trigger/non-trigger inputs and deterministic verification, filling a gap in instruction-following evaluation.
    \item We show that if-then capacity varies dramatically by instruction type (25--50 percentage point ranges) and that the difficulty ranking is consistent across models (Spearman $\rho = 0.83$, $p = 0.011$), establishing an inherent hierarchy of conditional instruction difficulty.
    \item We identify an asymmetry in failure modes: stronger models predominantly over-trigger (false positives) while weaker models predominantly miss triggers (false negatives), providing distinct design implications for practitioners.
    \item We demonstrate that the capacity bottleneck lies in the complexity of individual conditions rather than the number of simultaneous conditions, with both models maintaining near-perfect accuracy on up to 4 concurrent simple conditions.
\end{itemize}

The rest of this paper is organized as follows. \Secref{sec:related} reviews related work on instruction-following evaluation. \Secref{sec:method} describes the \benchname benchmark and evaluation protocol. \Secref{sec:results} presents our main results and analysis. \Secref{sec:discussion} discusses implications and limitations. \Secref{sec:conclusion} concludes.
